% !TEX program = xelatex
% 编译方式：xelatex 休学社保校招查证项目.tex（连续编译两次以生成正确的页码/超链接）
% 字体使用 TeX Live 自带的 Fandol 字体，Overleaf 与本地 TeX Live 均可直接编译
\documentclass[UTF8,fontset=fandol,zihao=-4]{ctexart}

\usepackage[a4paper,top=2.4cm,bottom=2.2cm,left=2.2cm,right=2.2cm]{geometry}
\usepackage[table]{xcolor}
\usepackage{enumitem}
\usepackage{booktabs}
\usepackage{longtable}
\usepackage{array}
\usepackage{needspace}
\usepackage{fancyhdr}
\usepackage{titlesec}
\usepackage[most]{tcolorbox}
\usepackage[colorlinks=true,urlcolor=link,linkcolor=link,pdftitle={休学期间缴纳社保对校招应届身份的影响：查证项目文档},pdfauthor={}]{hyperref}

% ---------- 配色 ----------
\definecolor{brand}{HTML}{1F4E79}
\definecolor{link}{HTML}{1565C0}
\definecolor{soft}{HTML}{EAF1F8}
\definecolor{note}{HTML}{FFF6E0}
\definecolor{noteline}{HTML}{E0A100}
\definecolor{quoteline}{HTML}{7A9CC6}
\definecolor{headbg}{HTML}{1F4E79}
\definecolor{rowalt}{HTML}{F4F7FB}
\definecolor{lowc}{HTML}{2E7D32}
\definecolor{midc}{HTML}{B26A00}
\definecolor{highc}{HTML}{C62828}

% ---------- 版式 ----------
\linespread{1.25}
\setlength{\parskip}{0.35em}
\setlength{\parindent}{0pt}
\setlist{itemsep=0.15em,topsep=0.3em,parsep=0pt,leftmargin=1.6em}
\setlist[enumerate]{leftmargin=1.8em}

\titleformat{\section}[block]{\zihao{3}\heiti\bfseries\color{brand}}{\chinese{section}、}{0pt}{}[{\vspace{-0.2em}\color{brand}\titlerule[0.8pt]}]
\titlespacing*{\section}{0pt}{1.6em plus 0.3em}{0.7em}
\titleformat{\subsection}[block]{\zihao{4}\heiti\bfseries\color{brand}}{}{0pt}{}
\titlespacing*{\subsection}{0pt}{1.1em plus 0.2em}{0.4em}

% 页眉页脚
\setlength{\headheight}{24pt}
\pagestyle{fancy}
\fancyhf{}
\renewcommand{\headrulewidth}{0.4pt}
\renewcommand{\headrule}{\hbox to\headwidth{\color{quoteline}\leaders\hrule height \headrulewidth\hfill}}
\fancyhead[L]{\small\color{gray}休学期间缴纳社保对校招应届身份的影响}
\fancyhead[R]{\small\color{gray}查证项目文档}
\fancyfoot[C]{\small\color{gray}第 \thepage\ 页}

% ---------- 自定义盒子 ----------
% 提示框
\newtcolorbox{notebox}{
  enhanced, breakable, colback=note, colframe=note,
  borderline west={2.5pt}{0pt}{noteline},
  boxrule=0pt, arc=2pt, left=8pt, right=8pt, top=5pt, bottom=5pt,
  before skip=6pt, after skip=8pt
}
% 引用/模板框
\newtcolorbox{quotebox}{
  enhanced, breakable, colback=soft, colframe=soft,
  borderline west={2.5pt}{0pt}{quoteline},
  boxrule=0pt, arc=2pt, left=8pt, right=8pt, top=5pt, bottom=5pt,
  before skip=4pt, after skip=8pt
}

% 加粗小标题
\newcommand{\lead}[1]{\par\Needspace{6\baselineskip}\vspace{0.3em}\noindent{\heiti\bfseries #1}\par}

% 表格辅助
\newcommand{\hc}[1]{\textcolor{white}{\heiti\bfseries #1}}
\renewcommand{\arraystretch}{1.35}
\newcolumntype{P}[1]{>{\raggedright\arraybackslash}p{#1}}
\newcommand{\lowrisk}{\textcolor{lowc}{\bfseries 低}}
\newcommand{\midrisk}{\textcolor{midc}{\bfseries 中}}
\newcommand{\midhigh}{\textcolor{midc}{\bfseries 中高}}
\newcommand{\highrisk}{\textcolor{highc}{\bfseries 高}}
\newcommand{\reqq}{\textcolor{highc}{\bfseries 必问}}
\newcommand{\optt}{\textcolor{gray}{选问}}
\newcommand{\todo}{\textcolor{highc}{（待填）}}

\begin{document}

\begin{tcolorbox}[enhanced, colback=brand, colframe=brand, arc=3pt,
  left=14pt, right=14pt, top=12pt, bottom=12pt, boxrule=0pt]
  {\color{white}
  {\zihao{2}\heiti\bfseries 休学期间缴纳社保对校招应届身份的影响}\\[0.5em]
  {\zihao{3}\heiti 查证项目文档（含完整提问清单与话术）}\\[0.8em]
  {\small 整理日期：2026 年 11 月 4 日（11 月 5 日更新：四地人社热线已咨询完毕） \textbar{} 用途：个人查证与行动清单 \textbar{} 状态：进行中}}
\end{tcolorbox}

\begin{notebox}
\small \textbf{说明：}本文的分析为一般性判断，不是定论。各单位、各年度的校招要求不同，\textbf{最终以招聘单位、学校和当地主管部门的书面或官方答复为准}。电话、网址在使用前请再核对一次官方渠道，因为这类信息可能变更。
\end{notebox}

\section{背景与上下文}

\lead{1. 基本情况}

\begin{itemize}
  \item 户籍：江苏
  \item 就读：大连的一所高校，本科
  \item 大四时因身体原因办理休学，休学期限两年
  \item 休学期间：在广东佛山工作，缴纳社保约 8 个月
  \item 后续计划：复学，以本科应届毕业生身份参加校招，目标单位为铁路局（国铁集团下属各铁路局）和中国中车
\end{itemize}

\lead{2. 起因}

向老师咨询“现在的工作能否交社保，会不会影响应届生身份”，老师表示问题较专业，且各单位要求不一，建议咨询人社局。因此启动本次查证。

\lead{3. 核心疑问}

\begin{enumerate}
  \item 休学期间缴纳过社保，毕业后还算不算校招意义上的“应届毕业生”？
  \item 这段工作和社保记录，会不会在铁路局、中车的审核、背调、体检入职阶段带来麻烦？
  \item 会不会影响毕业后的就业协议、档案、户口等手续？
\end{enumerate}

\lead{4. 待补充的关键信息（请你填写，会直接影响判断）}

\begingroup
\rowcolors{2}{rowalt}{white}
\begin{longtable}{@{}P{5.33cm}P{10.65cm}@{}}
\rowcolor{headbg}\hc{信息项} & \hc{你的情况}\\
\endfirsthead
\rowcolor{headbg}\hc{信息项} & \hc{你的情况}\\
\endhead
休学起止时间 & \todo{}\\
预计复学时间 & \todo{}\\
预计毕业年份 & \todo{}\\
社保目前是否已停缴 & \todo{}\\
佛山工作的劳动合同、离职证明是否齐全 & \todo{}\\
学校是否出具了休学、复学的正式文件 & \todo{}\\
目标岗位具体是哪几家铁路局、哪些中车子公司 & \todo{}\\
\end{longtable}
\endgroup

\lead{5. 最新进展：四地人社热线咨询结果（2026 年 11 月 5 日）}

已拨打佛山、常州、大连、南京四地人社热线（12333），完整记录见 4.8 咨询记录表。要点如下：

\begin{itemize}
  \item \textbf{大连、南京}：工作人员表示不了解此类认定，建议进一步咨询教育局和招录单位（如考公、进国央企，直接问招录单位）。
  \item \textbf{佛山}：首答不知道，转接司法相关人员后的口头答复是“休学期间的社保缴纳不影响应届生身份的判定”（无书面依据，仅供参考）。
  \item \textbf{常州}：答复“会影响”，追问后表示不敢确定，建议进一步咨询教育局。
\end{itemize}

\begin{notebox}
\small \textbf{阶段结论：}人社热线不负责“应届生身份”的认定，佛山与常州两个相反的口头答复也印证各地、各场景口径不一。人社这条线到此为止，不再拨打；重心转向\textbf{招录单位与教育部门}（见行动方法论新排序）。
\end{notebox}

\section{分析与判断}

\lead{1. 总体判断}

校招是否认可“应届”，通常看的是\textbf{学籍状态与毕业时间}，而不是有没有缴过社保。风险大概率不在“能不能报名”，而在“招聘单位如何看待这段工作经历”以及“毕业后档案、户口等手续如何衔接”。

\lead{2. 分点说明}

\begin{itemize}
  \item \textbf{学籍层面（风险较低）}：休学属于学籍管理中的正式状态，休学期间学籍仍保留在学校。只要你复学后正常毕业，你的应届届别一般按毕业年份认定。休学期间实习、兼职、工作并不少见。
  \item \textbf{校招资格层面（需向单位确认）}：铁路局、中车的校招公告一般写“当年度毕业的普通高校毕业生”，审核重点是学信网学籍、毕业证、毕业时间。但个别单位会在报名表或背调中询问“是否有工作经历”“是否缴纳过社保”，这属于需要\textbf{如实填写}的部分。
  \item \textbf{背景审查层面（需要留意）}：国企在政审、背调、入职材料核验阶段，可能看到你的社保缴纳记录。如实说明，并能用休学证明和劳动合同证明“这发生在休学期间”，通常是关键。
  \item \textbf{档案、户口层面（可能有差异）}：个别地方在认定应届身份用于就业登记、落户时，会把已缴社保视为已就业。这与校招资格不是一回事，但可能影响毕业后的档案接收、就业协议等手续，需向学校和就业部门确认。
  \item \textbf{地方政策参考（外地口径，仅作参照）}：上海规定地方国企和在沪央企校招不限制工作经历和社保，以毕业证书落款年度认定应届身份；重庆亦明确社保缴纳情况不影响报考。说明这类问题通常由教育、组织、国资部门联合规定，而非人社部门单独认定——与本次四地人社热线的答复情况一致。
\end{itemize}

\lead{3. 风险分级}

\begingroup
\rowcolors{2}{rowalt}{white}
\begin{longtable}{@{}P{3.79cm}P{2.52cm}P{4.21cm}P{4.63cm}@{}}
\rowcolor{headbg}\hc{风险点} & \hc{风险程度（初判）} & \hc{原因} & \hc{应对}\\
\endfirsthead
\rowcolor{headbg}\hc{风险点} & \hc{风险程度（初判）} & \hc{原因} & \hc{应对}\\
\endhead
能否报名校招 & \lowrisk{} & 通常看学籍和毕业时间 & 向招聘单位书面确认\\
审核时被问到工作经历 & \midrisk{} & 部分单位会问社保与工作经历 & 如实填写，备齐证明材料\\
被认定为“已就业”影响档案、户口 & \midrisk{} & 部分地区口径不同 & 问学校与教育部门（人社热线已确认不负责此认定）\\
社保仍在缴导致状态冲突 & \midhigh{} & 毕业前后状态不清晰 & 尽早确认并办理停保\\
隐瞒或表述不一致 & \highrisk{} & 背调发现会影响诚信评价 & 全程如实、材料一致\\
\end{longtable}
\endgroup

\lead{4. 一条原则}

不要靠猜，不要靠网上的二手经验。\textbf{让招聘单位和学校给出明确答复，并保留记录}，这比任何分析都有效。

\section{行动方法论}

\subsection{步骤 0：先整理材料（本周内可以完成）}

\begin{itemize}
  \item 休学申请与批准文件、复学申请与批准文件
  \item 佛山工作的劳动合同、离职证明、工资流水（可选）
  \item 社保缴纳记录（个人参保证明、缴费明细）
  \item 学信网学籍在线验证报告
  \item 一份时间线说明：休学 $\rightarrow$ 入职 $\rightarrow$ 离职 $\rightarrow$ 停保 $\rightarrow$ 复学
\end{itemize}

\lead{查询入口：}

\begin{itemize}
  \item 学信网：\url{https://www.chsi.com.cn}（查询学籍状态，下载在线验证报告）
  \item 社保记录：可通过“国家社会保险公共服务平台”或“广东政务服务网/粤省事”小程序查询缴费记录与参保证明
\end{itemize}

\subsection{步骤 1：问招聘单位（优先级最高）}

他们的答复最有效力。具体话术见第四部分 4.1。

\begin{itemize}
  \item \textbf{铁路局}：到你想报考的铁路局官网查看“招聘公告”，公告上通常留有咨询电话或邮箱。国铁集团系统的招聘信息也可通过各铁路局官网转到的招聘平台查看。
  \item \textbf{中国中车}：中国中车集团官网 \url{https://www.crrcgc.cc}，在“招聘”或“人才”栏目找到当年校招公告和联系方式；如目标是具体子公司（如中车四方、中车株机等），请到该子公司官网查公告。
  \item 做法：优先用\textbf{邮件或留言}，保留书面记录；电话咨询则记录时间和接线人。
\end{itemize}

\begin{notebox}
\small \textbf{提示：}具体网址和电话每年可能调整，请以当年校招公告上的联系方式为准。
\end{notebox}

\begin{notebox}
\small \textbf{2026 年 11 月检索发现（公告条款差异，问之前先对照原文）：}国铁集团下属中国铁信的招聘公告基本条件写有“纳入国家统一招生计划、\textbf{在校期间为非在职}的应届高校毕业生”——这一表述可能被解释为“休学期间有工作即属在职”，遇到此类条款务必当面或书面问清如何解释。另：郑州局公告写“列入国家招生计划的全日制本科应届毕业生”；上海局公告说明上一年已毕业离校未就业的毕业生可以应届身份应聘；中车部分子公司公告仅写“应届毕业生”、未展开定义。各单位口径不同，均以目标单位当年公告原文为准，对着条款逐条问。
\end{notebox}

\subsection{步骤 2：问学校（优先级高）}

话术见 4.2。

\begin{itemize}
  \item 联系对象：学院辅导员、学院教务办（学籍管理）、学校就业指导中心
  \item 要确认的核心问题：休学期间的工作和社保是否影响学籍、复学和毕业；就业协议、档案、户口迁移是否受影响；学校能否出具休学期间学籍保留的证明
  \item \textbf{大连交通大学}：教务处教务科 \textbf{0411-84106826}（沙河口校区 2 号教学楼 228 室）；旅顺口校区教学管理科 0411-86226605；招生与就业处就业办公室 \textbf{0411-84109958}（沙河口校区 2 号楼 209 室）；就业信息网 \url{https://dljt.bysjy.com.cn}；学校总机 0411-84106969
\end{itemize}

\begin{notebox}
\small \textbf{提醒：}以上电话为 2026 年 11 月网络检索结果，拨打前请以学校官网最新公布为准。教务科与就业办口径若不一致，请两边都留记录。
\end{notebox}

\subsection{步骤 3：问教育主管部门（优先级高，11 月 5 日新增）}

大连、南京两地人社热线均建议改问教育部门和招录单位，故新增此步骤。你在大连上学、学籍归辽宁，建议按以下顺序联系：

\begin{itemize}
  \item \textbf{辽宁省教育厅}：信访咨询电话 \textbf{024-86853225}，政务公开电话 024-86896570，官网 \url{https://jyt.ln.gov.cn}；接通后请其转“高校学生处／高校毕业生就业”对口部门
  \item \textbf{辽宁省大学生就业创业中心}（省教育厅下属，承办高校毕业生就业与派遣）：毕业生派遣咨询 \textbf{024-26901805}，辽宁毕业生服务网 \url{https://www.lnbys.com.cn}
  \item \textbf{大连市教育局}：综合／招考咨询 \textbf{0411-84649441}；师范类毕业生档案 0411-84638717
  \item \textbf{南京市教育局}：咨询电话 \textbf{025-83639901}，总机 025-83639900，官网 \url{https://edu.nanjing.gov.cn}
  \item 问法：先用 4.3 主模板报清事实，再对照追问清单；重点问清“休学期间缴纳社保是否影响应届身份认定、认定权在哪个部门、依据哪个文件”
\end{itemize}

\begin{notebox}
\small \textbf{提醒：}以上电话为 2026 年 11 月网络检索结果，拨打前请经官网核实。若教育局表示不归其管，追问“应由哪个部门负责解释”，并记录答复人和时间。
\end{notebox}

\subsection{步骤 4：问高校毕业生就业服务部门（优先级中）}

话术见 4.3。

\begin{itemize}
  \item \textbf{大连（学校所在地）}：高校毕业生就业与档案接收，一般由学校就业部门对接，必要时可咨询大连市人社部门，拨打 0411-12333（2026-11-05 已试：答复不了解应届认定，建议问教育局与招录单位）
  \item \textbf{全国大学生就业服务平台“24365”}：\url{https://24365.smartedu.cn}，可查看就业政策和咨询入口
  \item \textbf{江苏（户籍地）}：若涉及户口迁移、档案转递，可咨询户籍所在地人社部门，拨打 12333（需加所在城市区号，如南京 025-12333；应届认定问题人社热线已确认无法解答）
\end{itemize}

\subsection{步骤 5：问人社局（优先级较低，仅用于社保本身）}

人社局主要解决“社保记录和停保”的问题，\textbf{不负责“应届生”的认定}。话术见 4.4。

\begin{itemize}
  \item \textbf{2026 年 11 月 5 日已拨打佛山、常州、大连、南京四地 12333}，确认人社热线不给应届身份认定的权威答复（记录见 4.8），这条线不再重复拨打
  \item \textbf{佛山（仅剩社保业务）}：佛山市人社部门咨询电话 0757-12333；也可通过“佛山人社”公众号或“粤省事”小程序办理、咨询，用途限于停保确认、参保证明打印、跨省转移
  \item 12333 是全国统一的人社服务热线，人工服务通常集中在工作日白天
\end{itemize}

\subsection{步骤 6：兜底渠道}

\begin{itemize}
  \item \textbf{12345 政务服务便民热线}：12333 无法接通或问题不明确归属时使用，会转办到对应部门
  \item \textbf{国家政务服务平台}：\url{https://gjzwfw.www.gov.cn}，部分业务提供咨询或留言入口
\end{itemize}

\subsection{建议时间安排}

\begingroup
\rowcolors{2}{rowalt}{white}
\begin{longtable}{@{}P{2.66cm}P{9.42cm}P{3.48cm}@{}}
\rowcolor{headbg}\hc{阶段} & \hc{事项} & \hc{建议完成时间}\\
\endfirsthead
\rowcolor{headbg}\hc{阶段} & \hc{事项} & \hc{建议完成时间}\\
\endhead
第 1 周 & 整理材料；确认社保是否已停缴；查学信网学籍；向原单位 HR 索取证明（见 4.5） & 本周\\
第 1–2 周 & 电话大连交通大学教务处、就业办；联系辽宁省教育厅、大连市教育局、南京市教育局（见步骤 2、3） & 一周内\\
第 2–3 周 & 向目标铁路局、中车发邮件咨询，重点核对“在校期间为非在职”类条款（见步骤 1） & 两周内\\
第 3–4 周 & 视需要补问：常州人社到底指哪种场景的“影响”（落户、就业登记还是补贴）、依据哪个文件；必要时 12345 转办 & 视情况\\
校招季前 & 汇总所有答复，形成结论，准备报名材料 & 报名前\\
\end{longtable}
\endgroup

\section{提问清单与话术}

这一部分把可能需要问的问题按对象全部列出。\textbf{用法：}

\begin{enumerate}
  \item 先发各对象的“主模板”（邮件或留言），一次性把背景说清楚。
  \item 再对照“追问清单”补问：\reqq{}尽量一次问全，\optt{}视对方答复和你的实际情况再问。
  \item 拿到口头答复后，用 4.7 的“确认回复模板”发一条书面确认，把口头变成记录。
  \item 每次咨询都填入 4.8 的“咨询记录表”。
\end{enumerate}

\begin{notebox}
\small \textbf{提问的两条原则：}① 只陈述事实（学籍、休学、社保、时间），不要先下“我是应届生”的结论，让对方按他们的口径判断；② 涉及休学原因（健康方面），被问到时如实说明，没被问到不必主动展开病情细节，可以先问清楚“体检标准”再判断。
\end{notebox}

\needspace{7\baselineskip}
\subsection{4.1 给招聘单位（铁路局、中车）}

\lead{主模板（邮件或留言）}

\begin{quotebox}
\textbf{主题：}关于休学后复学的应届毕业生资格咨询（XX大学 XX 届 XXX）

您好！我是XX大学XX专业XX届本科生，因身体原因休学两年，休学期间（XXXX年X月—XXXX年X月）曾在佛山工作并缴纳社保约8个月，已于XXXX年X月离职并停保。我已（或将于XXXX年X月）复学，目前学信网学籍状态正常，预计XXXX年X月毕业，拟报考贵单位XXXX年校园招聘。

想请教几个问题：（1）我这种情况是否按应届毕业生对待、符合校招资格？（2）报名、审核、背调或入职阶段是否需要提供额外证明材料？（3）上述答复能否以邮件形式回复，便于我留存？

我可提供休学与复学证明、劳动合同、离职证明、社保缴费记录和学信网学籍验证报告。感谢您的帮助！

XXX｜电话：XXXXXXXXXXX｜邮箱：XXXX
\end{quotebox}

\lead{追问清单}

\begingroup
\rowcolors{2}{rowalt}{white}
\begin{longtable}{@{}P{0.87cm}P{7.18cm}P{5.43cm}P{1.65cm}@{}}
\rowcolor{headbg}\hc{\#} & \hc{问题} & \hc{为什么要问} & \hc{级别}\\
\endfirsthead
\rowcolor{headbg}\hc{\#} & \hc{问题} & \hc{为什么要问} & \hc{级别}\\
\endhead
1 & 休学后复学、以当年毕业身份报名，是否按应届毕业生对待？ & 核心结论 & \reqq{}\\
2 & 公告里“应届毕业生”的界定标准是什么？（毕业时间区间、学籍、是否要求未就业登记或未缴社保） & 对照公告原文，避免理解偏差 & \reqq{}\\
3 & 休学期间有社保缴纳记录，是否算“有工作经历”或“非应届”？是否影响资格？ & 直击你的情况 & \reqq{}\\
4 & 报名表、简历中是否要填工作经历？休学期间的经历应该怎么填？ & 保证表述与口径一致 & \reqq{}\\
5 & 审核、背调、政审会核查社保记录吗？在哪个阶段、核查哪些内容？ & 提前知道会被看到什么 & \reqq{}\\
6 & 需要提供哪些证明材料（休学证明、复学证明、离职证明、劳动合同）？有无盖章或格式要求？ & 提前备齐，避免临时补材料 & \reqq{}\\
7 & 是否要求入职前取得毕业证、学位证？报到截止日期是什么？ & 你的毕业时间可能晚于同级同学 & \reqq{}\\
8 & 答复能否以邮件或书面形式给我？ & 留存证据 & \reqq{}\\
9 & 对毕业年限、年龄是否有限制？休学导致年龄偏大是否有影响？ & 部分岗位设有年龄上限 & \optt{}\\
10 & 如果当年未能如期毕业，次年是否还能按应届报名？ & 预留退路 & \optt{}\\
11 & 入职体检的标准是什么？是否要求提供既往病史材料？目标岗位有无特殊身体条件要求？ & 与休学原因相关，提前对照 & \optt{}\\
12 & 就业协议书（三方协议）、档案、党团关系的接收和转递流程是什么？ & 衔接毕业后的手续 & \optt{}\\
13 & 江苏户籍，在贵单位所在地入职有无户籍或落户方面的要求和政策？ & 你的户籍在外省 & \optt{}\\
14 & 入职后社保、公积金如何衔接？佛山已缴的 8 个月是否需要转移，会不会影响试用期或缴费年限计算？ & 避免重复参保或断档 & \optt{}\\
15 & 如果审核时被认定为非应届，有没有补充材料或申诉的渠道？ & 提前了解兜底方案 & \optt{}\\
16 & 往年有没有类似情况（休学后复学、休学期间有工作经历）的录用先例？ & 参考他们的实际口径 & \optt{}\\
\end{longtable}
\endgroup

\lead{电话版开场白（约 30 秒）}

\begin{quotebox}
您好，我想咨询校园招聘的资格问题，请问现在方便说话吗？我是XX大学XX届的本科生，中间因病休学过两年，休学期间在外工作缴过约8个月社保，现在学籍正常、预计XXXX年毕业。想问一下，这种情况算不算贵单位认可的应届毕业生？报名和审核时需要注意什么？可以的话，能否请您把答复发到我邮箱，或者告诉我应该联系哪位老师？
\end{quotebox}

\subsection{4.2 给学校（辅导员、教务、就业中心）}

\lead{主模板（微信、邮件或当面均可）}

\begin{quotebox}
老师您好，我是XX学院XX专业XX级的XXX（学号XXXXXXXX）。我因身体原因于XXXX年X月休学，休学期间在外地工作并缴纳社保约8个月，现已离职，计划XXXX年X月复学。

想向您确认几件事：（1）休学期间的工作和社保记录，是否影响我的学籍、复学和毕业？（2）复学需要提交哪些材料、办理时间节点是什么？（3）毕业后就业协议书、档案、户口的办理会不会受影响？（4）学校能否为我出具休学期间学籍保留、预计毕业时间等证明？

耽误您时间了，谢谢老师！
\end{quotebox}

\lead{追问清单}

\begingroup
\rowcolors{2}{rowalt}{white}
\begin{longtable}{@{}P{0.87cm}P{7.18cm}P{5.43cm}P{1.65cm}@{}}
\rowcolor{headbg}\hc{\#} & \hc{问题} & \hc{为什么要问} & \hc{级别}\\
\endfirsthead
\rowcolor{headbg}\hc{\#} & \hc{问题} & \hc{为什么要问} & \hc{级别}\\
\endhead
1 & 休学期间在外工作、缴纳社保，是否违反学校的休学管理规定？ & 确认没有违规风险 & \reqq{}\\
2 & 复学资格如何审核？需要提交哪些材料（复学申请、医院诊断或康复证明、体检等）？时间节点是什么？ & 避免错过复学窗口 & \reqq{}\\
3 & 复学后编入哪个年级、预计哪一年毕业？学分和课程如何衔接？ & 确定应届届别 & \reqq{}\\
4 & 含休学在内的最长学习年限是多少？我会不会超限？ & 防止超过学制上限 & \reqq{}\\
5 & 休学和复学记录在学信网上如何显示？毕业年份会显示为哪一年？ & 招聘单位看的就是学信网 & \reqq{}\\
6 & 我在学校就业系统里的毕业届别是什么？能否正常使用就业推荐表、就业协议书？ & 确认就业手续不受影响 & \reqq{}\\
7 & 签就业协议书的条件是什么？已有社保记录会不会影响签署？ & 直击你的情况 & \reqq{}\\
8 & 我的档案现在存放在哪里？休学期间有无变动？毕业后档案接收、派遣流程是什么？会不会因有社保记录被认定为“已就业”而不能派遣？ & 档案是最容易被忽视的环节 & \reqq{}\\
9 & 学校能否出具：休学期间学籍保留证明、在读证明、复学证明、预计毕业证明？办理流程和时长？ & 后续向招聘单位提交 & \reqq{}\\
10 & 我的户口目前在学校集体户还是在江苏原籍？复学和毕业后的迁移规则是什么？ & 户口处理方式 & \optt{}\\
11 & 复学后学生医保（居民医保）怎么办理？之前在佛山参加的职工医保是否冲突、重复？ & 避免医保重复参保 & \optt{}\\
12 & 招聘单位来函做背景调查时，学校的统一答复口径和盖章流程是什么？ & 提前了解学校会怎么回答 & \optt{}\\
13 & 学校有没有类似情况（休学期间就业）的往年案例，后续是如何认定和办理的？ & 参考实际做法 & \optt{}\\
14 & 复学后能否继续参加校内招聘会、就业指导，学费和奖助学金是否受影响？ & 权益确认 & \optt{}\\
\end{longtable}
\endgroup

\subsection{4.3 给教育主管部门与就业服务部门（辽宁省教育厅、大连／南京市教育局、24365 平台）}

\lead{主模板（电话或留言）}

\begin{quotebox}
您好，我是大连XX大学的一名本科生，户籍在江苏XX市。我因病休学，休学期间在外地工作并缴纳社保约8个月，现已离职，计划复学并于XXXX年毕业。想咨询高校毕业生就业和档案方面的政策：（1）贵地对“应届高校毕业生”如何认定？（2）休学期间缴纳社保，会不会被视为已就业？（3）这对我毕业后的档案、就业协议和落户有没有影响？（4）这件事的认定权在哪个部门、依据哪个文件？
\end{quotebox}

\lead{追问清单}

\begingroup
\rowcolors{2}{rowalt}{white}
\begin{longtable}{@{}P{0.87cm}P{7.18cm}P{5.43cm}P{1.65cm}@{}}
\rowcolor{headbg}\hc{\#} & \hc{问题} & \hc{为什么要问} & \hc{级别}\\
\endfirsthead
\rowcolor{headbg}\hc{\#} & \hc{问题} & \hc{为什么要问} & \hc{级别}\\
\endhead
1 & “应届毕业生”在贵地的认定标准是什么？（毕业年限、学籍、就业登记状态） & 各地口径不同 & \reqq{}\\
2 & 休学期间缴纳的社保，是否会被视为“已就业”或“已就业登记”？ & 你最担心的点 & \reqq{}\\
3 & 如果被视为已就业，对毕业后档案接收、就业协议、落户、就业补贴、失业登记分别有什么影响？ & 评估后果大小 & \reqq{}\\
4 & 毕业后档案应转到哪里？江苏户籍是回原籍，还是迁到入职地？ & 衔接档案转递 & \reqq{}\\
5 & 毕业后离校未就业期间，有没有“择业期”？两年内是否仍按应届对待？ & 各地有不同规定 & \optt{}\\
6 & 如需办理就业登记或就业创业证，流程是什么？休学期间的社保是否影响办理？ & 后续手续 & \optt{}\\
7 & 对于休学期间有工作经历的学生，有没有专门的认定方式或证明要求？ & 找出是否有特例规定 & \optt{}\\
8 & 我应该向哪个具体部门或窗口提交材料或继续咨询？ & 防止被推诿 & \optt{}\\
\end{longtable}
\endgroup

\subsection{4.4 给人社部门（佛山）}

\lead{主模板（电话、公众号留言或小程序咨询）}

\begin{quotebox}
您好，我曾在佛山XX公司工作，参保时间为XXXX年X月至XXXX年X月（约8个月），现已离职。我想咨询社保的停缴、记录和后续衔接问题：（1）我的社保目前是否已停缴？（2）如何打印参保证明和缴费明细？（3）将来换到其他省市参保，这 8 个月如何转移和累计？
\end{quotebox}

\lead{追问清单}

\begingroup
\rowcolors{2}{rowalt}{white}
\begin{longtable}{@{}P{0.87cm}P{7.18cm}P{5.43cm}P{1.65cm}@{}}
\rowcolor{headbg}\hc{\#} & \hc{问题} & \hc{为什么要问} & \hc{级别}\\
\endfirsthead
\rowcolor{headbg}\hc{\#} & \hc{问题} & \hc{为什么要问} & \hc{级别}\\
\endhead
1 & 我的社保当前状态是什么？是否仍在缴费，或已显示中断？ & 确认没有继续被缴 & \reqq{}\\
2 & 原单位是否已办理减员？如果没有，我应该如何处理？ & 避免一直挂靠 & \reqq{}\\
3 & 如何打印参保证明和历史缴费明细（含各险种）？有没有电子凭证？ & 用作证明材料 & \reqq{}\\
4 & 在校复学后如果继续有社保记录，会有什么后果？如何确认已完全停止？ & 防止状态冲突 & \reqq{}\\
5 & 这 8 个月的养老保险缴费年限是否保留？将来跨省转移、接续怎么办理？ & 保护自己的权益 & \optt{}\\
6 & 医疗保险个人账户余额如何处理？转移或提取的规则是什么？ & 钱不要白放着 & \optt{}\\
7 & 复学后参加学生医保（居民医保）是否与职工医保冲突？ & 避免重复参保 & \optt{}\\
8 & 招聘单位或学校做背调时，能查到我社保记录的哪些内容？ & 预判会被看到什么 & \optt{}\\
9 & 原单位没有办停保或没有开离职证明，应如何投诉或维权？ & 遇到问题时的备选路径 & \optt{}\\
10 & 工伤、失业保险有没有需要处理的事项？ & 确认没有遗留事项 & \optt{}\\
\end{longtable}
\endgroup

\subsection{4.5 给原佛山单位 HR（索取证明材料）}

\lead{主模板}

\begin{quotebox}
XX 您好，我是 XXXX年X月至XXXX年X月在公司担任XX岗位的XXX。因后续需要提供在职和参保证明，想麻烦您协助：（1）开具离职证明（注明入职和离职日期、岗位）；（2）确认社保减员日期，并提供缴费明细；（3）如方便，提供劳动合同复印件。谢谢您的帮助！
\end{quotebox}

\lead{需要索取或确认的清单}

\begin{itemize}
  \item 离职证明（日期、岗位、离职原因表述中性、无争议）
  \item 劳动合同（电子版或复印件）
  \item 社保缴费明细及减员办理日期，确认没有继续缴纳
  \item 工资流水（可选）
  \item 公积金缴存情况及销户、转移方式（如有缴纳）
  \item 确认公司在档案、背调时的在职证明口径，保持与你的表述一致
\end{itemize}

\subsection{4.6 预料中的回答与应对}

\begingroup
\rowcolors{2}{rowalt}{white}
\begin{longtable}{@{}P{4.33cm}P{4.33cm}P{6.89cm}@{}}
\rowcolor{headbg}\hc{对方可能的回答} & \hc{说明什么} & \hc{你的下一步}\\
\endfirsthead
\rowcolor{headbg}\hc{对方可能的回答} & \hc{说明什么} & \hc{你的下一步}\\
\endhead
“以学籍和毕业证为准，社保不影响。” & 口径对你有利 & 请对方以邮件确认，保存截图或邮件，并记入咨询记录表。\\
“有社保记录就算有工作经历、不算应届。” & 对你不利，需要弄清依据 & 追问依据（公告第几条、哪个文件）；询问休学期间的情况是否有例外或可提交证明的补救方式；同时向学校和就业部门核实该口径。\\
“没有明确规定，建议问人社局。” & 对方也不确定 & 追问由谁负责解释；请对方告知具体部门；再向招聘负责人、学校就业中心交叉确认；不要只凭这一句下结论。\\
“不清楚，需要请示领导。” & 需要等待 & 约定回复时间（如三个工作日）；留下联系方式；到期礼貌催复。\\
“具体到背调时再看。” & 对方暂不判断 & 提前备齐证明材料；写好一页纸的情况说明；在报名表中如实填写。\\
“各地口径不一样，我们不好统一说。” & 对方不愿给结论 & 请对方指出以哪个部门的认定为准；改问学校就业中心和当地就业部门。\\
只给口头答复，不愿书面回复 & 缺少证据 & 你自己整理成文字，用确认邮件发给对方（见 4.7），请其回复“确认”或指出有误。\\
电话占线、无人接听 & 渠道问题 & 换时间段再打；改用邮件、公众号留言；必要时用 12345 转办。\\
长时间未回复 & 可能被遗漏 & 三个工作日后发催复；再不回复就换渠道或换联系人。\\
不同渠道答复互相矛盾 & 口径不统一 & 以招聘单位的书面答复为准；矛盾处再次书面请求对方说明；留存全部记录。\\
\end{longtable}
\endgroup

\subsection{4.7 跟进与确认话术}

\lead{确认回复模板（口头答复后发出）}

\begin{quotebox}
感谢您今天（X月X日）电话/当面的答复。我的理解是：XXXXXXXX（把对方的结论用一两句话复述）。如有理解不准确的地方，烦请指正；如无问题，也请您回复“确认”，方便我留存记录。谢谢！
\end{quotebox}

\lead{催复模板（三个工作日后）}

\begin{quotebox}
您好，我在X月X日发送了关于XXXX的咨询，想请问目前是否有进展？如需补充材料或说明，我可以随时提供。谢谢！
\end{quotebox}

\lead{感谢与结案模板}

\begin{quotebox}
感谢您的答复，我已按您的说明准备材料。后续如有变化，我会再向您确认。祝工作顺利！
\end{quotebox}

\subsection{4.8 咨询记录表}

每次咨询都记录：时间、渠道、对象和姓名（或工号）、原话要点、是否有书面答复、下一步。以下为示例行和空白行，可直接复制使用。

\begingroup
\rowcolors{2}{rowalt}{white}
\begin{longtable}{@{}P{1.59cm}P{1.34cm}P{2.26cm}P{2.76cm}P{3.01cm}P{1.09cm}P{1.84cm}@{}}
\rowcolor{headbg}\hc{日期} & \hc{渠道} & \hc{对象 / 部门} & \hc{我问的问题} & \hc{答复要点} & \hc{书面？} & \hc{下一步}\\
\endfirsthead
\rowcolor{headbg}\hc{日期} & \hc{渠道} & \hc{对象 / 部门} & \hc{我问的问题} & \hc{答复要点} & \hc{书面？} & \hc{下一步}\\
\endhead
例：11/10 & 邮件 & XX铁路局人力资源部 & 休学后能否按应届报名 & 以学籍和毕业时间为准 & 是 & 保存邮件\\
11/5 & 电话 & 大连人社热线 0411-12333 & 休学期间缴社保是否影响应届身份 & 不了解，建议咨询教育局与招录单位 & 否 & 转教育部门\\
11/5 & 电话 & 南京人社热线 025-12333 & 同上 & 不了解，建议咨询教育局与招录单位 & 否 & 转教育部门\\
11/5 & 电话 & 佛山人社热线 0757-12333 & 同上 & 首答不知道；转接司法人员后口头答复“休学期间的社保缴纳不影响应届生身份判定” & 否 & 仅口头，暂不作依据\\
11/5 & 电话 & 常州人社热线 0519-12333 & 同上 & 答复“会影响”；追问后表示不敢确定，建议咨询教育局 & 否 & 作罢，以教育部门与招录单位答复为准\\
\rule{0pt}{2.4em} &  &  &  &  &  & \\
\rule{0pt}{2.4em} &  &  &  &  &  & \\
\rule{0pt}{2.4em} &  &  &  &  &  & \\
\rule{0pt}{2.4em} &  &  &  &  &  & \\
\end{longtable}
\endgroup

\section{额外建议（小巧思）}

\begin{enumerate}
  \item \textbf{统一口径}：所有渠道、所有表格里，对这段经历的描述（时间、单位、岗位、是否在休学期间）要完全一致，避免背调时出现出入。
  \item \textbf{主动备好说明}：写一份一页纸的“休学期间工作经历情况说明”，附上休学证明、劳动合同和社保记录，被问到时可以直接提供。
  \item \textbf{同时问两家，对比答复}：铁路局和中车的口径可能不同，不要只问一家就下结论。
  \item \textbf{关注当年校招公告原文}：判断是否限制“有工作经历者”“已缴社保者”，要以公告条款为准。
  \item \textbf{健康状况相关}：校招体检和入职体检对部分岗位有要求，铁路系统对部分岗位的身体条件要求较明确，建议在报名前对照岗位的体检标准，提前了解自己是否符合。
  \item \textbf{医保衔接}：复学后注意学生医保和之前职工医保的衔接，避免重复参保或出现断档，具体向人社部门和学校确认（见 4.2、4.4）。
\end{enumerate}

\section{相关单位联系方式汇总}

{\small
\begingroup
\rowcolors{2}{rowalt}{white}
\begin{longtable}{@{}P{3.40cm}P{2.99cm}P{5.36cm}P{3.40cm}@{}}
\rowcolor{headbg}\hc{单位 / 渠道} & \hc{作用} & \hc{联系方式} & \hc{备注}\\
\endfirsthead
\rowcolor{headbg}\hc{单位 / 渠道} & \hc{作用} & \hc{联系方式} & \hc{备注}\\
\endhead
目标铁路局\todo{} & 确认校招资格、审核要求 & 官网“招聘公告”页的电话或邮箱 & 优先级最高，建议书面咨询\\
中国中车集团 / 目标子公司 & 确认校招资格、审核要求 & \url{https://www.crrcgc.cc}，以当年校招公告为准 & 优先级最高\\
大连交通大学 教务处（教务科） & 学籍、复学、毕业手续 & \textbf{0411-84106826}；旅顺口校区 0411-86226605 & 优先级最高（下一步就打）\\
大连交通大学 招生与就业处（就业办） & 就业协议、档案、户口、届别认定 & \textbf{0411-84109958}；就业信息网 \url{https://dljt.bysjy.com.cn} & 优先级最高（下一步就打）\\
辽宁省教育厅 & 学籍与毕业生就业政策口径 & 024-86853225（信访咨询）；024-86896570（政务公开） & 请其转高校学生处／就业对口部门\\
辽宁省大学生就业创业中心 & 毕业生派遣、就业手续 & 024-26901805；\url{https://www.lnbys.com.cn} & 省教育厅下属\\
大连市教育局 & 教育系统口径、招考综合 & 0411-84649441 & 2026-11 检索，打前核实\\
南京市教育局 & 教育系统口径 & 025-83639901（咨询）；025-83639900（总机） & 官网 edu.nanjing.gov.cn\\
本校学院辅导员 & 学籍、复学流程细节 & 学院公布方式\todo{} & 优先级高\\
学信网 & 查学籍状态、下载验证报告 & \url{https://www.chsi.com.cn} & 自助查询\\
24365 大学生就业服务平台 & 就业政策与咨询 & \url{https://24365.smartedu.cn} & 辅助\\
大连人社热线 & 社保、档案（不管应届认定） & 0411-12333 & 11/5 已试，不再问应届认定\\
南京人社热线 & 社保、档案（不管应届认定） & 025-12333 & 11/5 已试，不再问应届认定\\
常州人社热线 & 户籍地社保、落户（不管应届认定） & 0519-12333 & 11/5 已试；可补问“影响”所指场景\\
佛山人社热线 & 社保记录、停保 & 0757-12333；“佛山人社”公众号；“粤省事”小程序 & 优先级较低，主要查社保\\
原佛山单位 HR\todo{} & 离职证明、减员、缴费明细 & 公司 HR 联系方式 & 第 1 周内索取\\
全国人社服务热线 & 通用人社咨询 & 12333 & 加所在地区号拨打\\
政务服务便民热线 & 兜底、转办 & 12345 & 回复可能稍慢\\
国家政务服务平台 & 在线咨询、留言 & \url{https://gjzwfw.www.gov.cn} & 兜底\\
\end{longtable}
\endgroup
}

\begin{notebox}
\small \textbf{核对提醒：}表中的热线和网址请在使用前，通过各单位官方网站再确认一次；2026 年 11 月 5 日新增的电话均来自网络检索，尤其注意核实；标注“待填”的项目需要你补充具体信息。
\end{notebox}

\end{document}
